\documentclass[12pt]{article}
\usepackage[utf8]{inputenc}
\usepackage[T1]{fontenc}
\usepackage[margin=1in]{geometry}
\usepackage{setspace}
\usepackage{xcolor}
\usepackage[colorlinks=true, urlcolor=blue!60!black, linkcolor=blue!60!black]{hyperref}
\setlength{\parindent}{2em}
\setlength{\parskip}{1.25em}
\renewcommand{\baselinestretch}{1.0}
\begin{document}
\begin{center}
{\LARGE \textit{United Public Workers v. Mitchell}, 330 U.S. 75 (1947)\par}
\vspace{0.8em}
{\large Case Brief\par}
\end{center}
\vspace{1em}

\section{Memory Jogger}

In \href{https://www.courtlistener.com/opinion/104374/united-public-workers-of-america-v-mitchell/}{\textit{United Public Workers v. Mitchell}}, a 4--3 Supreme Court dismissed eleven federal employees' pre-enforcement challenges as nonjusticiable and upheld the Hatch Act's ban on partisan political management as applied to a U.S. Mint worker, declaring that Ninth and Tenth Amendment unenumerated-rights objections fail whenever Congress exercises a granted constitutional power.

\section{Detailed Case Facts}

In 1939, Congress enacted the Hatch Political Activity Act (53 Stat. 1147), amended in 1940 (54 Stat. 767), to curb partisan political activity by federal executive-branch employees and prevent the federal civil service from being mobilized into a partisan political machine (\href{https://www.law.cornell.edu/supremecourt/text/330/75}{330 U.S. 75 (1947)}; \href{https://firstamendment.mtsu.edu/article/united-public-workers-of-america-v-mitchell/}{First Amendment Encyclopedia}). The second sentence of \S 9(a) of the Act (18 U.S.C. \S 61h(a)) prohibited officers and employees in the executive branch---with exceptions for the President, Vice President, White House staff, department heads, and Senate-confirmed policymaking officers---from taking ``any active part in political management or in political campaigns,'' while preserving their right to vote and to express opinions on political subjects and candidates. Section 9(b) mandated immediate removal from office for any violation, and \S 15 (18 U.S.C. \S 61o) incorporated prior U.S. Civil Service Commission determinations defining prohibited partisan activities under Civil Service Rule I (\href{https://supreme.justia.com/cases/federal/us/330/75/}{Justia}).

Twelve individual federal executive-branch employees and their labor union, the United Public Workers of America (C.I.O.) (formerly the United Federal Workers of America), sued the members of the U.S. Civil Service Commission (headed by Harry B. Mitchell) seeking a declaratory judgment that the second sentence of \S 9(a) was unconstitutional and an injunction barring its enforcement (\href{https://www.law.cornell.edu/supremecourt/text/330/75}{Cornell LII}). The plaintiffs asserted that \S 9(a) violated: (1) the First Amendment freedoms of speech, press, and assembly; (2) the fundamental right of the people to engage in political activity reserved by the Ninth and Tenth Amendments; and (3) the Fifth Amendment Due Process Clause by arbitrarily depriving federal employees of liberty, discriminating between covered and exempted employees, and imposing vague prohibitions.

Crucially, the twelve individual plaintiffs stood in two very different factual postures:
\begin{itemize}
\item \textbf{Eleven Pre-Enforcement Plaintiffs:} Eleven employees (ranging from statisticians and economists to an arsenal lens grinder, a machinist, and a Navy Yard electric welder) had \emph{not} violated the Hatch Act. Their affidavits stated only that they desired, outside working hours, to write political articles, organize rallies, ring doorbells, distribute campaign literature, and act as poll watchers, but refrained out of fear of mandatory dismissal (\href{https://supreme.justia.com/cases/federal/us/330/75/}{Justia}).
\item \textbf{George P. Poole:} Only one plaintiff, George P. Poole---an industrial worker employed as a roller in the United States Mint in Philadelphia, Pennsylvania---had actually violated the Act. Outside working hours, Poole served as a Democratic Ward Executive Committeeman in Philadelphia's 28th Ward, worked at the polls on general election day (November 5, 1940), and acted as a paymaster distributing funds to other party workers. On January 12, 1944, the Civil Service Commission formally charged Poole with violating Civil Service Rule I and issued a proposed order for his removal unless he could refute the charges, which Poole instead admitted (\href{https://www.law.cornell.edu/supremecourt/text/330/75}{Cornell LII}).
\end{itemize}

\section{Procedural History}

Plaintiffs filed suit for injunctive and declaratory relief before a statutory three-judge panel of the U.S. District Court for the District of Columbia under 28 U.S.C. \S 380a and the Declaratory Judgment Act (28 U.S.C. \S 400). On September 26, 1944, the District Court held that all individual plaintiffs presented a justiciable controversy, but upheld \S 9(a) of the Hatch Act on the merits, dismissed the complaint, and granted summary judgment for the Commission (\textit{United Federal Workers of America (C.I.O.) v. Mitchell}, 56 F. Supp. 621 (D.D.C. 1944)) (\href{https://supreme.justia.com/cases/federal/us/330/75/}{Justia}).

The District Court entered an order allowing a direct appeal to the Supreme Court on October 26, 1944, under \S 3 of the Act of August 24, 1937 (28 U.S.C. \S 380a). Although \S 380a specified docketing within 60 days, the appeal was not docketed until February 2, 1945, prompting the Government to challenge the Supreme Court's appellate jurisdiction. The Supreme Court first heard oral argument on December 3, 1945, ordered reargument, heard reargument on October 17, 1946, and issued its decision affirming the District Court's judgment on February 10, 1947 (\href{https://tile.loc.gov/storage-services/service/ll/usrep/usrep330/usrep330075/usrep330075.pdf}{330 U.S. 75 (Library of Congress)}).

\section{Judicial Votes}

The case was decided by a \textbf{4--3} vote on the merits (affirming dismissal as to Poole), with \textbf{two Justices not participating} (\href{https://www.law.cornell.edu/supremecourt/text/330/75}{Cornell LII}; \href{https://supreme.justia.com/cases/federal/us/330/75/}{Justia}):
\begin{itemize}
\item \textbf{Majority Opinion (4 Justices on the merits):} Justice Stanley F. Reed delivered the opinion of the Court, joined by Chief Justice Fred M. Vinson, Justice Harold H. Burton, and Justice Felix Frankfurter (who also wrote separately on jurisdiction).
\item \textbf{Concurring Opinion:} Justice Felix Frankfurter (arguing the appeal should have been dismissed as untimely docketed under 28 U.S.C. \S 380a, but reaching the merits and joining Justice Reed's opinion under compulsion of the Court's jurisdictional ruling).
\item \textbf{Dissenting Opinion:} Justice Hugo L. Black (voting that all twelve plaintiffs presented a justiciable controversy and that \S 9(a) of the Hatch Act was facially unconstitutional under the First and Fifth Amendments).
\item \textbf{Opinion Dissenting in Part:} Justice William O. Douglas (disagreeing with the majority on two of the four issues: arguing that all eleven other plaintiffs presented a justiciable controversy, and arguing on the merits that \S 9(a) was unconstitutional as applied to an industrial worker like Poole).
\item \textbf{Dissenting Statement:} Justice Wiley B. Rutledge (dissented as to Poole for the reasons stated by Justice Black, and declined to pass on the constitutional claims of the other eleven appellants because he agreed the controversy as to them was not yet appropriate for discretionary declaratory judgment jurisdiction---making the vote \textbf{5--2} against entertaining the eleven pre-enforcement claims).
\item \textbf{Not Participating:} Justice Frank Murphy and Justice Robert H. Jackson took no part in the consideration or decision of the case.
\end{itemize}

\section{Holding}

\textbf{Judgment:} Affirmed (56 F. Supp. 621, affirmed).

\textbf{Rule Announced:}
\begin{enumerate}
\item \textbf{Appellate Jurisdiction:} The 60-day docketing provision of 28 U.S.C. \S 380a is not a jurisdictional bar; absent a motion to dismiss by the appellee under Supreme Court Rules 11 and 47, the Court retains power to hear the appeal.
\item \textbf{Justiciability (Article III):} Federal employees who merely express a general desire to engage in prohibited political activities in the future, without having violated the statute or facing an actual threat of disciplinary action beyond the statute's existence, do not present a concrete ``case or controversy'' under Article III. By contrast, an employee who admits violating the Act and faces a formal proposed removal order (Poole) presents a justiciable controversy suitable for declaratory relief.
\item \textbf{Substantive Constitutionality (First, Fifth, Ninth, and Tenth Amendments):} Congress may, within reasonable limits, prohibit executive-branch employees---including industrial workers acting off-duty---from taking an active part in partisan political management or political campaigns in order to promote efficiency, discipline, and integrity in the federal civil service. Moreover, when a challenger invokes unenumerated rights reserved by the Ninth and Tenth Amendments against federal action, the inquiry focuses on whether Congress possesses a granted power; ``[i]f granted power is found, necessarily the objection of invasion of those rights, reserved by the Ninth and Tenth Amendments, must fail'' (\href{https://www.law.cornell.edu/supremecourt/text/330/75}{330 U.S. at 96}; \href{https://constitution.congress.gov/browse/essay/amdt9-3/ALDE_00013643/}{Constitution Annotated}).
\end{enumerate}

\textbf{Holding vs. Dicta:}
\begin{itemize}
\item \textbf{Holdings:} (1) Late docketing under 28 U.S.C. \S 380a did not strip the Court of jurisdiction; (2) the eleven non-violating employees' pre-enforcement claims were nonjusticiable for lack of a concrete controversy; (3) Poole's claim was justiciable; and (4) \S 9(a) of the Hatch Act and Civil Service Rule I were constitutional as applied to Poole's off-duty service as a ward executive committeeman, poll worker, and party paymaster, notwithstanding the First, Fifth, Ninth, and Tenth Amendments.
\item \textbf{Dicta:} Justice Reed's sweeping concession that Congress could not ``enact a regulation providing that no Republican, Jew or Negro shall be appointed to federal office, or that no federal employee shall attend Mass or take any active part in missionary work'' (330 U.S. at 100) was illustrative dicta acknowledging outer Bill of Rights limits on public employment conditions. Likewise, the Court expressly declined to examine the outer reaches of \S 15's incorporation of Civil Service Commission rules beyond Poole's clear violations, reserving facial vagueness and overbreadth questions (which the Court later addressed and rejected in \href{https://www.law.cornell.edu/supremecourt/text/413/548}{\textit{U.S. Civil Service Comm'n v. National Ass'n of Letter Carriers}}, 413 U.S. 548 (1973)).
\end{itemize}

\section{Analysis of Opinions}

\subsection{Majority Opinion (Justice Reed)}

Justice Reed addressed four sequential issues. First, he held that the 60-day docketing timeframe in 28 U.S.C. \S 380a was a procedural rule governed by Supreme Court Rules 11 and 47 rather than a strict jurisdictional limit (\href{https://supreme.justia.com/cases/federal/us/330/75/}{Justia}). Second, invoking Article III's ban on advisory opinions, Reed concluded that the eleven employees who had not violated the Hatch Act presented only an ``abstract'' attack on the Act's political expediency; courts cannot adjudicate constitutional limits on hypothetical future speech without ``definite rights'' on one side and ``definite prejudicial interferences'' on the other. Third, Reed held that Poole presented a concrete, justiciable controversy because the Commission had formally charged him and issued a proposed removal order, and removal was mandatory upon his admitted facts.

Fourth, reaching the merits solely as to Poole, Reed framed the constitutional inquiry as a deferential balance between individual political liberty and Congress's authority to maintain an efficient, nonpartisan civil service. Crucially for substantive due process and unenumerated-rights doctrine, Reed addressed the plaintiffs' Fifth, Ninth, and Tenth Amendment claims by treating the Ninth and Tenth Amendments as mere residual mirrors of enumerated federal powers: because federal powers are ``subtracted from the totality of sovereignty originally in the states and the people,'' once a granted congressional power is found, any Ninth or Tenth Amendment objection ``must fail'' (\href{https://constitution.congress.gov/browse/essay/amdt9-3/ALDE_00013643/}{Constitution Annotated}). Relying on \textit{Ex parte Curtis}, 106 U.S. 371 (1882), and \textit{United States v. Wurzbach}, 280 U.S. 396 (1930) (which upheld bans on political financial contributions by and to federal employees), Reed reasoned that the Hatch Act merely extended those principles from ``political contributions of money'' to ``political contributions of energy.'' Because the Act left untouched employees' right to vote and to express non-partisan opinions on public affairs, Congress could rationally conclude that even an off-duty industrial worker like Poole might be swept into a political machine or have his promotion tied to partisan loyalty. Whether administrative and industrial workers should be treated differently was a ``matter of detail for Congress,'' and courts should interfere ``only when such regulation passes beyond the generally existing conception of governmental power'' (\href{https://www.law.cornell.edu/supremecourt/text/330/75}{Cornell LII}).

\subsection{Concurring Opinion (Justice Frankfurter)}

Justice Frankfurter disagreed solely on the threshold jurisdictional issue. Reading \S 3 of the Act of August 24, 1937 (28 U.S.C. \S 380a) in light of Congress's intent to expedite review of constitutional challenges to federal statutes, Frankfurter argued that Congress exercised its Article III power to regulate appellate jurisdiction by imposing a mandatory 60-day docketing deadline that the Court had no power to enlarge by rule (\href{https://www.law.cornell.edu/supremecourt/text/330/75}{Cornell LII}). However, ``under compulsion of the Court's assumption of jurisdiction,'' Frankfurter reached the merits and joined Justice Reed's opinion in full.

\subsection{Dissenting Opinion (Justice Black)}

Justice Black, agreeing with Justice Douglas that all twelve plaintiffs presented a justiciable controversy, argued that \S 9(a) of the Hatch Act was unconstitutional on its face under the First and Fifth Amendments (\href{https://www.law.cornell.edu/supremecourt/text/330/75}{Cornell LII}). Black emphasized the staggering sweep of the statute, which muzzled roughly three million federal employees and thousands of federally funded state employees---and even chilled their spouses under Commission rules holding employees accountable for family members' political activity. The statute's purported preservation of the right to ``express their opinions'' was a ``hopeless contradiction'' with the Commission's sweeping bans, leaving employees to speak ``at their peril'' and reducing them to ``mere spectators'' in American democracy. Rejecting the majority's deference, Black insisted that laws restricting First Amendment freedoms must be narrowly drawn to address a ``grave and imminent danger'' rather than punishing millions of law-abiding citizens to deter coercion by a few supervisors. Congress could directly criminalize official coercion or patronage abuses without ``hack[ing] at the roots of a Government by the people themselves.''

\subsection{Opinion Dissenting in Part (Justice Douglas)}

Justice Douglas dissented on two grounds (\href{https://supreme.justia.com/cases/federal/us/330/75/}{Justia}). First, on justiciability, he argued that the eleven non-violating employees faced an immediate, non-hypothetical threat of mandatory discharge; forcing lower-income civil servants to sacrifice their livelihoods and seniority just to test their constitutional rights defeated the core purpose of the Declaratory Judgment Act. Second, on the merits, Douglas drew a sharp distinction---modeled on British civil service practice---between \emph{administrative} employees (whose proximity to policymaking and public contact might justify partisan neutrality rules) and \emph{industrial} workers like Poole (a skilled Mint roller as remote from policymaking ``as a charwoman''). Warning that expanding public employment could politically sterilize vast swaths of the electorate, Douglas declared that federal employees ``are not second-class citizens'' and that their fundamental political rights could not be sacrificed absent a ``clear and present danger to the civil service system.''

\subsection{Dissenting Statement (Justice Rutledge)}

Justice Rutledge issued a brief separate statement dissenting as to Poole for the reasons stated by Justice Black, while declining to pass on the constitutional claims of the other eleven appellants because he believed their claims were not yet appropriate for the discretionary exercise of declaratory judgment jurisdiction (\href{https://www.law.cornell.edu/supremecourt/text/330/75}{Cornell LII}).

\section{Examples: Future Applications}

\begin{enumerate}
\item \textbf{Same-Side (Upheld): Federal Statute Barring Executive-Branch Civil Servants from Running as Partisan Candidates for State Legislature.} Suppose Congress enforces a Hatch Act provision prohibiting a GS-11 federal customs inspector from running as a Democratic or Republican nominee for a state legislative seat, even if all campaigning occurs off-duty. Under \textit{Mitchell} (and its reaffirmation in \href{https://www.law.cornell.edu/supremecourt/text/413/548}{\textit{Letter Carriers}}, 413 U.S. 548 (1973)), the restriction is \textbf{constitutional}. Running as a partisan party nominee is core ``active part in political management or in political campaigns,'' and Congress's interest in preserving an impartial, merit-based civil service outweighs the employee's First, Fifth, and Ninth Amendment objections.
\item \textbf{Same-Side (Nonjusticiable): Pre-Enforcement Declaratory Challenge by a Federal Employee Alleging Only a General Desire to ``Engage in Political Organizing.''} Suppose a federal postal clerk files a federal declaratory judgment action alleging that she wishes to ``participate actively in unspecified future political campaigns'' and fears firing under the Hatch Act, without identifying any concrete planned conduct or receiving any agency warning or disciplinary notice. Under \textit{Mitchell}'s Article III holding, the suit is \textbf{nonjusticiable} because federal courts will not render advisory opinions on abstract, hypothetical disputes lacking definite threatened interference.
\item \textbf{Opposite-Side (Invalidated): Federal Regulation Barring Civil Servants from Attending Religious Services or Discriminating Against Members of a Political Party in Hiring.} Suppose an agency issues a rule barring any federal employee from attending Mass, engaging in off-duty religious missionary work, or holding membership in a particular mainstream political party. Even under \textit{Mitchell}'s deferential framework, this regulation is \textbf{unconstitutional}. Justice Reed explicitly recognized that Congress may not condition federal employment on surrendering religious exercise or bar citizens from office based on party affiliation, race, or religion (330 U.S. at 100); modern unconstitutional-conditions doctrine (\href{https://www.law.cornell.edu/supremecourt/text/427/347}{\textit{Elrod v. Burns}}, 427 U.S. 347 (1976)) reinforces that viewpoint or religious discrimination passes beyond legitimate governmental power.
\item \textbf{Opposite-Side (Invalidated): Blanket Ban on Federal Employees Receiving Compensation for Off-Duty, Non-Partisan Speeches and Articles.} Suppose Congress prohibits all executive-branch employees below Grade GS-16 from accepting any honoraria for off-duty speeches or essays on topics completely unrelated to their jobs or partisan politics (such as literary criticism or dance reviews). As the Supreme Court later held in \href{https://www.law.cornell.edu/supremecourt/text/513/454}{\textit{United States v. National Treasury Employees Union}}, 513 U.S. 454 (1995), this broad restriction is \textbf{unconstitutional}. Unlike the Hatch Act's targeted ban on partisan campaign management upheld in \textit{Mitchell}, an indiscriminate ban on non-partisan expressive activity imposes a sweeping burden on citizen speech unsupported by a demonstrated threat to workplace efficiency.
\item \textbf{Fence-Sitter: Disciplining a Low-Level Federal Mechanic for Hosting a Non-Partisan Issue Advocacy Blog That Frequently Criticizes the Sitting Administration.} Suppose a federal maintenance technician at a national park maintains an off-duty personal blog criticizing federal climate and labor policies and urging readers to vote against incumbents who support those policies, without formally endorsing any political party or working for a party committee. Under \textit{Mitchell}, the case sits squarely on the fence. On the one hand, Justice Reed emphasized that the Hatch Act leaves unrestricted ``[e]xpressions, public or private, on public affairs, personalities and matters of public interest, not an objective of party action'' (330 U.S. at 100), and modern public-employee speech doctrine (\href{https://www.law.cornell.edu/supremecourt/text/461/138}{\textit{Connick v. Myers}}, 461 U.S. 138 (1983)) protects off-duty citizen speech on matters of public concern. On the other hand, if the Commission characterizes coordinated electoral advocacy against specific incumbents as directing energy ``toward party success,'' \textit{Mitchell}'s deferential standard toward Congress and the Commission would weigh in favor of upholding discipline.
\end{enumerate}

\section{Critique}

\textbf{Originalist and Textualist Critiques (The Ninth Amendment and Unenumerated Rights):} Within the Substantive Due Process and unenumerated-rights canon, \textit{Mitchell} is infamous for Justice Reed's dismissive treatment of the Ninth Amendment. Originalist scholar Randy E. Barnett has sharply critiqued Reed's ``rights-powers conception''---the claim that whenever Congress possesses a granted power, any Ninth Amendment objection ``necessarily ... must fail''---arguing that it effectively erases the Ninth Amendment from the Constitution by reducing it to a redundant truism (\href{https://scholarship.law.cornell.edu/cgi/viewcontent.cgi?article=3396&context=clr}{Randy E. Barnett, \textit{Reconceiving the Ninth Amendment}, 74 Cornell L. Rev. 1 (1988)}). On Barnett's originalist view, the Ninth Amendment protects individual natural rights that affirmatively constrain how enumerated powers may be exercised, rather than merely describing whatever space remains outside federal power. From a competing originalist-federalism perspective, Professor Kurt T. Lash argued that Reed's New Deal-era pairing of the Ninth and Tenth Amendments as mere ``truisms'' broke with a nineteenth-century jurisprudence (rooted in James Madison and Justice Joseph Story) that treated the Ninth Amendment as a rule of construction restricting expansive readings of federal power to preserve local self-government (\href{https://scholarship.richmond.edu/context/law-faculty-publications/article/2469/viewcontent/The_Lost_Jurisprudence_of_the_Ninth_Amendment__2005.pdf}{Kurt T. Lash, \textit{The Lost Jurisprudence of the Ninth Amendment}, 83 Tex. L. Rev. 597 (2005)}). Nearly two decades after \textit{Mitchell}, the Court pivoted sharply away from Reed's Ninth Amendment minimalism in \textit{Griswold v. Connecticut}, 381 U.S. 479 (1965), invoking the Ninth Amendment alongside Bill of Rights penumbras to protect unenumerated marital privacy (\href{https://constitution.congress.gov/browse/essay/amdt9-3/ALDE_00013643/}{Constitution Annotated}).

\textbf{Progressive and Doctrinal Critiques (Unconstitutional Conditions and Public-Employee Speech):} Progressive and First Amendment scholars have faulted \textit{Mitchell} as a lingering relic of the discredited ``right-privilege distinction'' epitomized by Justice Holmes's 1892 epigram in \textit{McAuliffe v. Mayor of New Bedford} (``a constitutional right to talk politics, but ... no constitutional right to be a policeman''), which Reed cited favorably in Footnote 34 (\href{https://chicagounbound.uchicago.edu/cgi/viewcontent.cgi?article=4265&context=uclrev}{U. Chi. L. Rev.}; \href{https://scholarship.law.wm.edu/cgi/viewcontent.cgi?article=2634&context=wmlr}{Wm. \& Mary L. Rev.}). Commentators note that \textit{Mitchell} applied an unusually lax, near-rational-basis deference to severe restrictions on core political speech, failing to require narrow tailoring or less restrictive alternatives (such as distinguishing industrial workers from policymaking officials, as Justice Douglas urged). Although the Supreme Court reaffirmed \textit{Mitchell}'s bottom-line validation of the Hatch Act in \href{https://www.law.cornell.edu/supremecourt/text/413/548}{\textit{Letter Carriers}} (1973), later cases dismantled the right-privilege premise under the unconstitutional-conditions doctrine, and Congress itself eventually relaxed the Hatch Act's off-duty restrictions through the Federal Employees Political Activities Act of 1993 (\href{https://firstamendment.mtsu.edu/article/united-public-workers-of-america-v-mitchell/}{First Amendment Encyclopedia}).

\textbf{Underlying Competing Values:} \textit{Mitchell} pits two core democratic commitments against one another: (1) \textit{nonpartisan administrative integrity and protection against patronage coercion}---ensuring that an expanding federal bureaucracy serves all citizens impartially and does not degenerate into an entrenched one-party spoils machine; and (2) \textit{equal political citizenship and personal autonomy}---ensuring that millions of public employees retain full First Amendment expression and unenumerated personal liberty off the clock without being relegated to second-class citizenship as the price of government employment.

\section{Key Quotations}

\begin{quote}
``The powers granted by the Constitution to the Federal Government are subtracted from the totality of sovereignty originally in the states and the people. Therefore, when objection is made that the exercise of a federal power infringes upon rights reserved by the Ninth and Tenth Amendments, the inquiry must be directed toward the granted power under which the action of the Union was taken. If granted power is found, necessarily the objection of invasion of those rights, reserved by the Ninth and Tenth Amendments, must fail.''
\end{quote}
\noindent --- Justice Stanley F. Reed, opinion of the Court, 330 U.S. at 95--96 (\href{https://supreme.justia.com/cases/federal/us/330/75/}{Justia}; \href{https://www.courtlistener.com/opinion/104374/united-public-workers-of-america-v-mitchell/}{CourtListener}).

\begin{quote}
``The determination of the extent to which political activities of governmental employees shall be regulated lies primarily with Congress. Courts will interfere only when such regulation passes beyond the generally existing conception of governmental power.''
\end{quote}
\noindent --- Justice Stanley F. Reed, opinion of the Court, 330 U.S. at 102 (\href{https://supreme.justia.com/cases/federal/us/330/75/}{Justia}).

\begin{quote}
``The section of the Act here held valid reduces the constitutionally protected liberty of several million citizens to less than a shadow of its substance. It relegates millions of federal, state, and municipal employees to the role of mere spectators of events upon which hinge the safety and welfare of all the people, including public employees.''
\end{quote}
\noindent --- Justice Hugo L. Black, dissenting, 330 U.S. at 115 (\href{https://supreme.justia.com/cases/federal/us/330/75/}{Justia}).

\begin{quote}
``Absent coercion, improper use of government position or government funds, or neglect or inefficiency in the performance of duty, federal employees have the same rights as other citizens under the Constitution. They are not second-class citizens.''
\end{quote}
\noindent --- Justice William O. Douglas, dissenting in part, 330 U.S. at 124 (\href{https://supreme.justia.com/cases/federal/us/330/75/}{Justia}).

\begin{quote}
``In that Act, Congress put a limit to the time within which a case may be docketed here after an appeal below is allowed. Such a limitation by Congress is in the exercise of its power to regulate the appellate jurisdiction of this Court. It is not within our power to enlarge a limit fixed by Congress unless Congress itself gave the Court such dispensing power.''
\end{quote}
\noindent --- Justice Felix Frankfurter, concurring, 330 U.S. at 104 (\href{https://supreme.justia.com/cases/federal/us/330/75/}{Justia}).

\end{document}
